\section{Threats to Validity}
\label{sec:threats}

\textit{Internal validity.}
The completeness score formula (Equation~\ref{eq:completeness}) and its component weights
(25/25/20/15/15) were defined a priori based on domain judgment rather than optimized on
a holdout set. Different weight assignments would shift absolute scores but are unlikely to
change the relative ordering of diseases or the qualitative patterns we report.
To verify this, we tested three alternative weight configurations: equal weights
(20/20/20/20/20), a coverage-heavy scheme (35/35/10/10/10), and an accuracy-heavy scheme
(20/20/10/25/25). The disease ranking is preserved across all configurations: Dengue,
Measles, and COVID-19 consistently score highest and Ebola consistently lowest,
confirming that the completeness ordering is robust to reasonable weight variation.
The structural error taxonomy reflects common failure modes observed during development; it may
not cover all possible errors in other disease models.

\textit{Construct validity.}
Fuzzy string matching with synonym groups is a heuristic: it may produce false positives
when different epidemiological concepts share a word (e.g., ``Susceptible adults'' matching
``Susceptible children'') or false negatives for non-standard abbreviations not covered by
our synonym list. The 0.6 \texttt{SequenceMatcher} threshold was set to minimize both error
types on our dataset but was not tuned on a separate validation set.

\textit{External validity.}
Our study covers 10 diseases with models of varying complexity (4--23 flows). These were
chosen to span a range of standard structures (SIR, SEIR, multi-group, host--vector) but
may not represent the full diversity of published compartmental models, particularly agent-based
or network-based models that do not fit the ODE compartmental structure. Results for
diseases with highly non-standard parameter naming conventions (e.g., Bayesian notation)
may differ from what we report here.
Although evaluated on these 10 ODE compartmental models, the framework is
design-agnostic at the pipeline level: the three-phase structure (reference construction,
LLM-assisted extraction, gap completion) applies to any parametric model class where a
formal reference and a natural-language description coexist, including pharmacokinetic,
ecological, and within-host models.

\textit{LLM non-determinism.}
All Phase~2 extractions and Phase~3 inferences used temperature~0, which minimizes but
does not eliminate non-determinism in some LLM backends. A small fraction of results could
differ on replication. To mitigate this, we report averaged patterns and note that all runs
used the same prompts and model versions.

\textit{Gold standard quality.}
The reference \texttt{.compmodel} files were constructed from each paper's ODE equations
and parameter tables, and validated against the reported simulation results and figures
in the paper, not from informal reading alone.
Each compartment, flow, and parameter in the gold model traces directly to a named
equation or table entry. For papers that provided simulation code, we cross-validated
the gold model's compartment wiring and initial conditions against the code. For those
without code, ambiguous parameter assignments were resolved by returning to the
differential equation system and identifying which flow each symbol multiplies.
This construction methodology ensures the gold standard is grounded in the paper's own
formal specification, minimising subjective interpretation. The reference models and
their source annotations are included in the project repository.
Any residual misinterpretation of ambiguous notation propagates to the completeness score,
which is itself the instrument for detecting such ambiguity in the source paper.